\section{Study 2: independent targets and strong competitors}
\label{sec:study2}

Study~2 evaluated \STwoTargetsTotal\ independent targets (\STwoTargetsMin\ to
\STwoTargetsMax\ per benchmark). Table~\ref{tab:study2} lists every frozen decision with the
number of benchmarks on which a Holm-significant difference favours EOST-M or the competitor,
the number on which EOST-M is better, and the median relative improvement (for C3 the
relative difference of the PDM Bayesian CAT, for C5a the range of coverage; C5a, C5b and C6
concern the recalibrated interval, marked by an asterisk). Of the \STwoNHyp\ decisions,
\STwoNSupported\ were supported. Apart from C5b, every superiority contrast favoured EOST-M
on at least seven of the eight benchmarks and no competitor was significantly better on any
benchmark; four of the five superiority hypotheses that were not supported (C2, C4, C7, C8)
missed the threshold by one benchmark. CRUXEval-O, with only \STwoTargetsMin\ independent
clusters, reached significance in no contrast.

\begin{table}[htbp]
\centering\footnotesize
\caption{Study 2: frozen hypotheses at $N=8$ and their decisions.}
\label{tab:study2}
\setlength{\tabcolsep}{3pt}
\begin{tabular}{@{}lllcccl@{}}
\toprule
 & EOST-M versus & Endpoint & Holm sig. & positive & median & Supported \\
\midrule
R2 & plug-in 2PL (MFI:pi) & log loss & 5/0 & 8/8 & $+5.9$ & yes \\
R3 & Naive & score IAE & 6/0 & 8/8 & $+53.9$ & yes \\
R4 & plug-in 2PL (MFI:pi) & $|\mathrm{cov}-0.9|$ & 5/0 & 8/8 & $+41.6$ & yes \\
C1 & Rasch-type plug-in & log loss & 5/0 & 7/8 & $+2.4$ & yes \\
C2 & Rasch-type marginal & log loss & 4/0 & 8/8 & $+1.8$ & no \\
C4 & Random:mg & score IAE & 4/0 & 8/8 & $+21.1$ & no \\
C5b & raw EOST-M interval$^\ast$ & interval score & 0/2 & 4/8 & $+1.0$ & no \\
C6 & Naive (Wilson)$^\ast$ & interval score & 6/0 & 7/8 & $+35.0$ & yes \\
C7 & ATLAS-style 3PL CAT & score IAE & 4/0 & 8/8 & $+26.0$ & no \\
C8 & tinyBenchmarks & error @50 & 4/0 & 8/8 & $+29.6$ & no \\
C9 & MINCE & error @50 & 5/0 & 8/8 & $+44.8$ & yes \\
\midrule
C3 & PDM Bayesian CAT (equivalence) & log loss & \multicolumn{2}{c}{4/8 within $\pm$1\%} & $+0.32$ & no \\
C5a & (coverage in $[0.85,0.95]$)$^\ast$ & coverage & \multicolumn{2}{c}{5/8 banks inside} & 0.81--0.91 & no \\
\bottomrule
\end{tabular}
\end{table}


\paragraph{Replication with independent targets.} The study-1 results held: EOST-M's
held-out log loss was lower than plug-in Fisher CAT's on every benchmark (R2, median
\STwoRTwoMed\%, \STwoRTwoFav\ significant), more so with four calibrators
(\SecRTwoNFourMed\%) than with sixteen (\SecRTwoNSixteenMed\%), and by the same amount for
deterministic (\STwoRTwoDet\%) and aggregate verdicts (\STwoRTwoAgg\%); R3 and R4 were also
supported, and plug-in Fisher CAT's intervals covered the score only \STwoCovmfipi\ of the
time, against \STwoCoveigmg\ for EOST-M.

\paragraph{Integration, constraint and Bayesian CATs.} EOST-M predicted better than the
Rasch-type plug-in test (C1, median \STwoCOneMed\%), more so with sixteen calibrators
(\SecCOneNSixteenMed\%) than with four (\SecCOneNFourMed\%), where a single discrimination
becomes a visible misspecification; the integrated one-parameter model sat in between (C2,
\STwoCTwoMed\%). The two tests had similar score errors (median \SecEostVsRaschIAEMed\%, no
benchmark significant), and the Rasch-type intervals were again better calibrated before
recalibration (\STwoCovmfiOnepr). The PDM Bayesian CAT differed from EOST-M by a median of
\STwoCThreeMed\% in log loss and was within $\pm1\%$ on only \STwoCThreeEq\ benchmarks, so
equivalence (C3) was not supported. Its frozen burn-in proved too short on
\FbDriftFrozenBank, where the mean calibrator ability drifted by a median of
\FbDriftFrozenMax\ during the retained sweeps; with either burn-in (10 or, post hoc, 60
sweeps) it was not significantly different from EOST-M on any benchmark, in log loss or score
error (corrected log loss higher by a median of \FbBmiMed\% with Bayesian maximum information
and \FbEigMed\% with EIG). The exploratory joint fully Bayesian CAT, which
keeps the uncertainty about the location of the calibration shared across items, was the
closest to EOST-M with EIG selection (median \FbjEigMed\%). At $N=4$ the corrected PDM
Bayesian CAT was never significantly better than EOST-M (within $\pm1\%$ on \FbBmiEqNFour\
and \FbEigEqNFour\ benchmarks with the two criteria). With eight calibrators, integrating
their abilities and the hyperprior on top of the item parameters thus made no significant
difference, but equivalence within $\pm1\%$ was not shown.

\paragraph{Adaptivity and benchmark reduction.} All methods used the same eight calibrators;
ATLAS and tinyBenchmarks were designed for hundreds to thousands, so these comparisons concern
the small-panel regime, not the methods in their native settings. Under the same calibrated
model, adaptive selection reduced the score error of random items by a median of
\STwoCFourMed\% (C4, \STwoCFourFav\ benchmarks significant), more with larger panels
(\SecCFourNFourMed\% at $N=4$, \SecCFourNSixteenMed\% at $N=16$). Table~\ref{tab:reduction}
gives the score errors at 25, 50 and 100 items, averaged over targets and benchmarks;
tinyBenchmarks and MINCE choose their items without adapting to the target. After 50 items
EOST-M's error was \STwoCEightMed\% below tinyBenchmarks' gp-IRT (C8) and \STwoCNineMed\%
below MINCE's subset mean (C9, supported), and over the test \STwoCSevenMed\% below the
ATLAS-style CAT's (C7), whose plug-in 3PL parameters and randomesque selection made it the
least accurate model-based policy (coverage \STwoCovatlas). With eight calibrators,
tinyBenchmarks' K-Means cannot always find as many anchors as its budget, because its item
embeddings take only as many values as there are distinct response patterns; an audit
decided before the results were seen showed that restricting C8 to full-budget cells, or
comparing with EOST-M after the number of items tinyBenchmarks actually used, again gave
\TinyFullFav\ significant benchmarks. MINCE's own stopping rule asked for a median of
\STwoMinceNstarMed\ items and then had a score error of \STwoMinceErrNstar\ points, more
than EOST-M after 100 items (\STwoErreigmgHundred).

\begin{table}[htbp]
\centering\footnotesize
\caption{Study 2: absolute score error (points) after 25, 50 and 100 items at $N=8$.}
\label{tab:reduction}
\begin{tabular}{@{}lccc@{}}
\toprule
Method & 25 items & 50 items & 100 items \\
\midrule
EOST-M (EIG:mg) & 3.39 & 2.81 & 2.13 \\
ATLAS-style 3PL CAT & 4.46 & 4.20 & 3.79 \\
Random:mg & 4.78 & 3.34 & 2.42 \\
tinyBenchmarks gp-IRT & 4.96 & 4.39 & 4.49 \\
tinyBenchmarks p-IRT & 4.69 & 3.92 & 3.31 \\
tinyBenchmarks anchors & 6.23 & 5.62 & 5.37 \\
MINCE subset mean & 6.71 & 4.69 & 3.19 \\
Naive (Wilson) & 7.05 & 4.88 & 3.26 \\
\bottomrule
\end{tabular}
\end{table}


\paragraph{Calibrated uncertainty.} Model-based intervals remained too narrow: EOST-M
covered the score \STwoCoveigmg\ of the time, the PDM Bayesian CAT \FbBmiCov, the Rasch-type
test \STwoCovmfiOnepr\ and plug-in Fisher CAT \STwoCovmfipi, whereas Random:mg reached
\STwoCovrandommg\ and Naive's Wilson interval \STwoCovnaive. Recalibration raised EOST-M's
coverage to \STwoCoveigmgRR, inside $[0.85,0.95]$ on \STwoCFiveaIn\ benchmarks, one short of
the criterion (C5a), and did not improve the interval score (C5b; significantly worse on
HellaSwag and WinoGrande): the wider intervals cost about as much as the misses they removed.
The recalibrated intervals had a much better interval score than Naive's Wilson intervals
(C6, median \STwoCSixMed\%) because they were much narrower (mean width \STwoWidtheigmgRR\
against \STwoWidthnaive). Recalibration brought the Rasch-type test inside the band on all
eight benchmarks, but with intervals about twice as wide (\STwoWidthmfiOneprRR). It needs
calibrators as pseudo-targets and helped more as the panel grew (coverage
\STwoCovEostRNFour, \STwoCovEostRNEight\ and \STwoCovEostRNSixteen\ at $N=4$, 8 and 16), and
mixing one uniformly drawn item in four into the adaptive test raised coverage only to
\STwoCovmixmg.

\paragraph{Temporal validation and the hyperprior.} On the \TempTargets\ SWE-bench
Verified submissions of the latest quarter, calibrated only on earlier submissions,
\TempNSig\ of the \TempNTests\ frozen contrasts favoured EOST-M at a Holm-adjusted $p<0.05$
within the component and none favoured a competitor; EOST-M's log loss was \TempRTwoMed\% lower than plug-in
Fisher CAT's and \TempCOneMed\% lower than the Rasch-type test's. Because the empirical-Bayes
hyperprior is essential and estimated from the same panel, it was checked post hoc and varied
under the frozen plan. Resampling items moves the hyperparameters by at most
\EBItemsSdMaxEight\ (standard deviation, $N=8$), whereas resampling calibrators moves them
about as much as the variation between panels: the hyperprior reflects which systems are in
the panel, not noise in the items. Prior predictive checks reproduce the spread of the
items' proportions correct, but the share of items with identical panel verdicts only on
average, and real systems agree with each other less, item by item, than a unidimensional
2PL with independent items implies (mean item--rest correlation \PPCRestObs\ observed,
\PPCRestSim\ simulated): the response model is misspecified in its dependence structure, one
reason why model-based intervals are too narrow. A cross-fitted hyperprior, estimated for
each item on the other half of the items, left EOST-M's log loss unchanged (median
\SecXfVsEostMedAbs\%), and Student-$t$ hyperpriors made it slightly worse
(\SecTpVsEostMedAbs\%): the gain of EOST-M does not come from fitting the prior to the items
it is applied to.
